\section{Finite-Window Identifiability and Information Bounds}
\label{sec:identifiability}

\subsection{Information Structure}

For accepted scalar rows, define $h_C=q/C_b$ and $h_R=k_RI_\Sigma$, with conditional variances $R_C$ and $R_R$. Because the two observations act along disjoint parameter directions, the finite-window information Gramian is
\begin{equation}
\vect{G}_{\theta,N}=
\sum_{k\in\mathcal K_C}\frac{1}{R_{C,k}}
\begin{bmatrix}h_{C,k}^2&0\\0&0\end{bmatrix}
+
\sum_{k\in\mathcal K_R}\frac{1}{R_{R,k}}
\begin{bmatrix}0&0\\0&h_{R,k}^2\end{bmatrix}.
\label{eq:gramian}
\end{equation}

Suppose every accepted charge interval has length at least $T_{w,C,\min}$ and satisfies $|i_C|\ge I_{C,\min}^{\mathrm{eval}}$ without changing sign. Then $|q|\ge q_{\min}^{\mathrm{eval}}=I_{C,\min}^{\mathrm{eval}}T_{w,C,\min}$. If $R_C\le R_{C,\max}$, $R_R\le R_{R,\max}$, the calibrated $k_R$ remains fixed with known conservative evaluated-envelope lower bound $k_{R,\min}^{\mathrm{eval}}>0$, and $I_\Sigma\ge I_{\Sigma,\min}^{\mathrm{eval}}>0$, windows containing $m_C$ charge rows and $m_R$ edge rows have the lower bounds
\begin{align}
\underline\mu_C & =m_C\frac{(I_{C,\min}^{\mathrm{eval}}T_{w,C,\min}/C_b)^2}{R_{C,\max}},
\label{eq:mu_c}\\
\underline\mu_R & =m_R\frac{(k_{R,\min}^{\mathrm{eval}})^2(I_{\Sigma,\min}^{\mathrm{eval}})^2}{R_{R,\max}}.
\label{eq:mu_r}
\end{align}

\begin{proposition}[Operating-condition-induced finite-window identifiability]
\label{prop:identifiability}
Assume valid CCM; $k_R$ known and fixed during each identification window; sign-invariant accepted charge windows; positive accepted edge-current sum; finite conditional covariances; and at least one accepted charge and edge row in a finite window. Then
\begin{equation}
\vect{G}_{\theta,N}\succeq
\operatorname{diag}(\underline\mu_C,\underline\mu_R)\succ\vect{0},
\label{eq:positive_gramian}
\end{equation}
and $(C,r_C)$ is locally identifiable over the accepted finite window.
\end{proposition}

\emph{Proof:} Each charge row contributes a positive semidefinite matrix only in the $\alphab$ direction, and each edge row contributes one only in the $r_C$ direction. Applying the physical lower bounds in~\eqref{eq:mu_c} and~\eqref{eq:mu_r}, summing $m_C$ and $m_R$ accepted contributions, and noting that both lower bounds are strictly positive establishes~\eqref{eq:positive_gramian}. Since $C=C_b/\alphab$ is locally one-to-one for $\alphab>0$, local identifiability transfers from $(\alphab,r_C)$ to $(C,r_C)$. \hfill$\square$

The result is conditional: PWM alone does not guarantee persistent excitation at every load or conduction mode. The \Cuk{} topology supplies the physical directions, while accepted rows certify when the bound applies. The superscript ``eval'' marks minima over the evaluated CCM envelope; these constants are neither online gates nor post-hoc tuning values.

\subsection{Numerical Lower-Bound Check}

Across 45 independent Model-B CCM cases, the weakest physical quantities were
\begin{align}
I_{C,\min}^{\mathrm{eval}}&=0.330108~\mathrm{A},&
I_{\Sigma,\min}^{\mathrm{eval}}&=0.438504~\mathrm{A},\nonumber\\
q_{\min}^{\mathrm{eval}}&=0.660217~\mu\mathrm{C}.&&
\end{align}
with $T_{w,C,\min}=2.0~\mu$s and $k_{R,\min}^{\mathrm{eval}}=0.957654$. The resulting weakest information bounds were
\begin{equation}
\underline\mu_C=0.217943,\qquad
\underline\mu_R=880.008.
\label{eq:numerical_mu}
\end{equation}
The minimum ratios of empirical information to the corresponding conservative bound were 1.02872 and 1.02334. Hence every retained case exceeded its computed bound, while the weakest case rather than a nominal case determines the reported constants.

\begin{figure}[t]
\centering
\includegraphics[width=\columnwidth]{fig_pe_lower_bound.png}
\caption{Theoretical lower bounds versus empirical finite-window information over the Model-B CCM cases. Circles denote normalized inverse capacitance and squares denote physical ESR; their parameterization-dependent magnitudes and units are not directly comparable. All accepted cases lie above the equality line.}
\label{fig:pe_lower_bound}
\end{figure}

\subsection{Covariance Contraction for the Uncertainty Extension}

For the uncertainty extension, let $P_n$ be a posterior variance, $Q_N$ the process variance accumulated over a finite window, and $\mu\ge\underline\mu>0$ the accepted information. Then
\begin{equation}
P_{n+1}\le
\frac{P_n+Q_N}{1+\mu(P_n+Q_N)}.
\label{eq:cov_recursion}
\end{equation}

\begin{proposition}[Finite-window covariance boundedness]
\label{prop:covariance}
For recurring valid updates, finite $0\le Q_N\le\bar Q_N$, and $\mu\ge\underline\mu>0$, the posterior covariance satisfies~\eqref{eq:cov_recursion}. For constant $Q_N$, its positive fixed point is
\begin{equation}
P^*=\frac{-Q_N+\sqrt{Q_N^2+4Q_N/\mu}}{2}.
\label{eq:fixed_point}
\end{equation}
For $Q_N=0$, $P_n^{-1}\ge P_0^{-1}+n\mu$.
\end{proposition}

This conditional linear-model result is consistent with standard filtering theory \cite{jazwinski1970filtering,anderson1979optimal}; it is not a global nonlinear-converter convergence theorem. Appendix~\ref{app:covariance} gives the derivation and validated fixed-point values. Projection was inactive in 200 legal-PE seeds and is treated as a safety constraint rather than the source of contraction.
